\documentclass[letterpaper,12pt]{article}
%\usepackage{harvard}
%\usepackage{subfigure}
%\usepackage[active]{srcltx}
\usepackage{graphicx}
\usepackage{pdflscape}
\usepackage{amsmath}
\usepackage{amsthm}
\usepackage{lscape}
\usepackage{enumerate}
\usepackage{float}
\usepackage{grffile}
\usepackage{relsize}
\usepackage{tablefootnote}
\usepackage{footnote}
\usepackage{sectsty}
\usepackage{xcolor}
\usepackage{caption}
\usepackage{color}
\usepackage{natbib}
\usepackage{multirow}
\usepackage{bbm}
\usepackage[T1]{fontenc}
\usepackage[multiple]{footmisc}
\usepackage{appendix}
\usepackage{times}
\usepackage{relsize}
\usepackage{caption}
\usepackage{newtxtext,newtxmath}
\usepackage{enumitem}
\usepackage{booktabs}


%%%%%%%%%%%%%%%%%%%%%%%%%%%%%%%%%%%%%%%%%%%%%%%%%%%%%%%%%%%%
%%%   Title Page Information
%%%%%%%%%%%%%%%%%%%%%%%%%%%%%%%%%%%%%%%%%%%%%%%%%%%%%%%%%%%%
\title{ECON311: Assignment I}
\author{}
\date{Due date: TBD, Fall 2026}

%%%%%%%%%%%%%%%%%%%%%%%%%%%%%%%%%%%%%%%%%%%%%%%%%%%%%%%%%%%%
%%%   Double and Single Space Commands
%%%%%%%%%%%%%%%%%%%%%%%%%%%%%%%%%%%%%%%%%%%%%%%%%%%%%%%%%%%%
\newlength{\defbaselineskip}
\setlength{\defbaselineskip}{\baselineskip}
\newcommand{\setlinespacing}[1]{\setlength{\baselineskip}{#1 \defbaselineskip}}
\newcommand{\doublespacing}{\setlength{\baselineskip}{1.3 \defbaselineskip}}
\newcommand{\singlespacing}{\setlength{\baselineskip}{1\defbaselineskip}}
\renewcommand{\thesubsection}{\arabic{subsection}}
\newtheoremstyle{dotlessS}{}{}{\itshape}{}{\bfseries}{}{ }{}
\theoremstyle{dotlessS}
\newtheorem{theorem}{Theorem}
\newtheorem{lemma}{Lemma}
\newtheorem{proposition}{Proposition}
\newtheorem{assumption}[theorem]{Assumption}
\newtheorem{corollary}{Corollary}
\newtheorem{definition}{Definition}
\newenvironment{example}[1][Example]{\begin{trivlist}
\item[\hskip \labelsep {\bfseries #1}]}{\end{trivlist}}
\newenvironment{remark}[1][Remark]{\begin{trivlist}
\item[\hskip \labelsep {\bfseries #1}]}{\end{trivlist}}

\newcommand\Pitem{%
  \addtocounter{enumi}{-1}%
  \renewcommand\theenumi{\arabic{\enumi}'}%
  \item%
  \renewcommand\theenumi{\arabic{enumi}}%
}

%%%%%%%%%%%%%%%%%%%%%%%%%%%%%%%%%%%%%%%%%%%%%%%%%%%%%%%%%%%%
%%%   Set Margins Here
%%%%%%%%%%%%%%%%%%%%%%%%%%%%%%%%%%%%%%%%%%%%%%%%%%%%%%%%%%%%
\textwidth=6.5in
\textheight=9in
\oddsidemargin=0.10in
\evensidemargin=0.10in
\topmargin=-0.5in
%\parskip=3pt

%%%%%%%%%%%%%%%%%%%%%%%%%%%%%%%%%%%%%%%%%%%%%%%%%%%%%%%%%%%%%
%%%   Actual Document Begins Here
%%%%%%%%%%%%%%%%%%%%%%%%%%%%%%%%%%%%%%%%%%%%%%%%%%%%%%%%%%%%%
\begin{document}
\pagenumbering{arabic}
\maketitle
\vspace{-.3in}
\textbf{Instructions:} There are $5$ questions for a total of $80$ points. Answer all questions. To receive full marks, write your answers as complete and clear as possible.

%----------------------------------------------------------------------
\section*{Question I: National Income Accounting [10 points]}
%----------------------------------------------------------------------
The table below presents national income and expenditure data for the hypothetical economy of \textbf{Northland} in the year 2025. All figures are in billions of dollars.

\vspace{0.3em}
\begin{table}[H]
\centering
\begin{tabular}{lc}
\toprule
\textbf{Item} & \textbf{Amount (\$B)} \\
\midrule
Consumer spending on goods and services & 4{,}000 \\
Gross business fixed investment & 1{,}700 \\
Change in business inventories & 200 \\
Government purchases of goods and services & 2{,}200 \\
Government transfer payments (pensions, social assistance) & 500 \\
Exports & 1{,}000 \\
Imports & 700 \\
\midrule
Wages and salaries (labour income) & 5{,}000 \\
Corporate profits and rents (capital income) & 2{,}200 \\
Depreciation of the capital stock & 1{,}200 \\
\bottomrule
\end{tabular}
\end{table}

\begin{enumerate}[label=(\alph*)]
\item (4 points) Compute GDP for Northland in 2025 using the \textbf{expenditure approach}. Clearly identify each component ($C$, $I$, $G$, $NX$). Should government transfer payments be included in $G$? Explain why or why not. What is Northland's net investment (gross investment minus depreciation)?

\item (3 points) Using the \textbf{income approach}, verify that you obtain the same GDP as in part~(a). What are the components of GDP from the income side? In one or two sentences, explain why the expenditure approach and the income approach must give the same answer in theory.

\item (3 points) Real GDP in Northland in 2025 is $\$8{,}000$ billion (measured in constant 2024 prices). In 2024, nominal GDP equalled real GDP (i.e., 2024 is the base year, with $GDP^{deflator}_{2024} = 100$). Compute the GDP deflator for 2025 and the inflation rate between 2024 and 2025 as measured by the GDP deflator.
\end{enumerate}


%----------------------------------------------------------------------
\section*{Question II: Cross-Country Development}
%----------------------------------------------------------------------
For this question you will need to download four data series. You can find all data at the World Bank Development Indicators website: \\
\textcolor{blue}{https://databank.worldbank.org/source/world-development-indicators}. The series needed are: i.~``GDP per capita (constant 2015 $US\$$)'', ii.~``GDP per capita, PPP (constant 2021 international $\$$)''. Obtain these data for both \textbf{Canada} and \textbf{South Korea} from \textbf{1995 to 2022}. You are free to use any software; Microsoft Excel is recommended.

\begin{enumerate}[label=(\alph*)]
\item (10 points) Plot GDP per capita (constant 2015 $US\$$) for both countries on the same graph with time on the horizontal axis. On a separate graph, plot GDP per capita, PPP (constant 2021 international $\$$) for both countries. Which series better captures differences in living standards across countries? Explain the economic reason for the difference between the two series.

\item (10 points) Assume that in any given year the GDP of a country is given by $Y=\overline{A}(K)^{\frac{1}{3}}(L)^{\frac{2}{3}}$, and assume as in class that there are no TFP differences across countries ($\overline{A}$ is the same for both). Using the GDP per capita, PPP series, calculate the implied ratio of capital per worker in Canada relative to South Korea for every year from 1995 to 2022. Plot this ratio over time. Has the gap in capital per worker been \textbf{narrowing or widening} over this period? What does the Solow model predict should happen to this ratio over time if both countries share the same steady state? Is your finding consistent with that prediction? Explain.

\item (10 points) Using GDP per capita, PPP (constant 2021 international $\$$), calculate the average annual growth rates for both the Canadian and South Korean economies over the entire 1995--2022 period. Assume South Korea continues to grow at a constant rate equal to its 1995--2022 average. How many years from 2022 would it take for South Korea's GDP per capita to reach Canada's 2022 level?
\end{enumerate}


%----------------------------------------------------------------------
\section*{Question III: Growth Accounting [10 points]}
%----------------------------------------------------------------------
The table below shows the average annual growth rates of real GDP ($g_Y$), the capital stock ($g_K$), and the labour force ($g_L$) over a 20-year period for two hypothetical economies, Alpha and Beta.

\vspace{0.5em}
\begin{table}[H]
\centering
\begin{tabular}{lccc}
\toprule
\textbf{Economy} & $g_Y$ & $g_K$ & $g_L$ \\
\midrule
Alpha & 3.0\% & 4.5\% & 1.5\% \\
Beta  & 6.0\% & 6.0\% & 0.0\% \\
\bottomrule
\end{tabular}
\end{table}

\noindent Assume that both economies have the production function $Y = \overline{A}(K)^{\frac{1}{3}}(L)^{\frac{2}{3}}$ and use the growth-accounting identity:
\begin{eqnarray*}
g_Y \approx g_A + \tfrac{1}{3}\,g_K + \tfrac{2}{3}\,g_L
\end{eqnarray*}

\begin{enumerate}[label=(\alph*)]
\item (5 points) For each economy, compute: (i) the contribution of capital accumulation to GDP growth, (ii) the contribution of labour force growth to GDP growth, and (iii) the rate of TFP growth $g_A$. What fraction of total GDP growth is attributable to TFP growth in each economy? Which economy relies more on TFP as a source of growth, and what might explain this difference?

\item (5 points) Compute the growth rate of GDP per capita ($g_y$) for each economy. Using the per-capita growth-accounting identity $g_y \approx g_A + \frac{1}{3}(g_K - g_L)$, decompose GDP per capita growth into (i) the contribution from \textbf{capital deepening} (growth in capital per worker) and (ii) the contribution from \textbf{TFP growth}. Comment on the relative importance of these two sources for each economy's living standard improvement.
\end{enumerate}


%----------------------------------------------------------------------
\section*{Question IV: Technology and the Solow Model [15 points]}
%----------------------------------------------------------------------
Suppose that at time $t_0$, a country experiences a \textbf{permanent improvement in total factor productivity} — its TFP level jumps from $\overline{A}$ to $\overline{A}' > \overline{A}$ and remains at $\overline{A}'$ forever (for example, due to a technological breakthrough or the adoption of superior production methods). Assume the economy was at its initial steady state just before $t_0$.

Using the Solow model, explain what happens to the economy both \textbf{at the instant of the shock} ($t = t_0$) and \textbf{over time} afterward. Discuss the levels of capital ($K$), capital per worker ($k$), GDP ($Y$), and GDP per worker ($y$) both during the transition and at the new steady state. You are \textbf{not} required to discuss how growth rates change during the transition.

In your answer:
\begin{enumerate}[label=(\roman*)]
\item Draw the Solow diagram showing the effect of the TFP increase.
\item Draw separate time-path diagrams for $K$, $k$, $Y$, and $y$.
\item \textbf{Carefully address the following:} Does output per worker $y$ change \textit{immediately} at $t_0$, before any new capital has accumulated, or does it change only gradually over time? Why? How does this feature of a TFP shock differ from a change in the savings rate (studied in lectures)?
\end{enumerate}


%----------------------------------------------------------------------
\section*{Question V: Endogenous Growth with a Growing Labour Force}
%----------------------------------------------------------------------
Consider a Solow model in which the production function exhibits no diminishing returns to capital: $Y_t = \overline{A}\,K_t$. Unlike the basic version of this model discussed in class, assume that the labour force grows at a constant rate $n > 0$.

\begin{enumerate}[label=(\alph*)]
\item (5 points) Derive the capital-accumulation equation in per-worker terms (i.e., in terms of $k_t = K_t/L_t$). Draw the Solow diagram for this economy in per-worker terms with $k_t$ on the horizontal axis. Clearly label all curves and their slopes.

\item (5 points) Suppose the economy begins with capital per worker $k_0 > 0$. Describe and illustrate how the economy evolves over time for each of the following two cases: \textbf{(i)} $\overline{s}\,\overline{A} > \overline{d} + n$ and \textbf{(ii)} $\overline{s}\,\overline{A} < \overline{d} + n$. In each case, is there a stable steady state? What is the economic interpretation?

\item (5 points) Derive the condition under which GDP per capita grows at a strictly positive, constant rate in the long run. What is that constant growth rate? How does this long-run prediction differ from the prediction of the standard Solow model with $Y_t = \overline{A}\,K_t^{1/3}\,L^{2/3}$? Provide an intuitive explanation for why introducing population growth changes the condition for sustained per-capita growth in the AK model.
\end{enumerate}


\end{document}
